% =======> Project Final Report for CSE 151B Spring 2026 Math Reasoning Competition
% Authors: Connor Wu
% =================================================

\documentclass{article}

\PassOptionsToPackage{numbers, compress}{natbib}

\usepackage[preprint]{neurips_2024}

\usepackage[utf8]{inputenc}
\usepackage{amsmath}
\usepackage{amssymb}
\usepackage{mathtools}
\usepackage{amsthm}
\usepackage{float}
\DeclareMathOperator{\R}{\mathbb{R}}
\usepackage{microtype}
\usepackage{graphicx}
\usepackage{subcaption}
\usepackage{booktabs}
\usepackage{algorithm}
\usepackage{algorithmic}
\usepackage{hyperref}
\usepackage{wrapfig}
\usepackage[capitalize,noabbrev]{cleveref}
\RequirePackage{doi}
\usepackage{xcolor, colortbl}
\definecolor{darkblue}{rgb}{0.19, 0.19, 0.62}
\hypersetup{
  colorlinks = true,
  citecolor  = darkblue,
  filecolor  = darkblue,
  urlcolor   = darkblue,
}
\usepackage[textsize=tiny]{todonotes}


\title{
CSE 151B/251B Final Report:\\
Inference-Time Engineering for Math Reasoning\\
with a Thinking Language Model
}

\author{%
  Connor Wu \\
  Department of Computer Science and Engineering\\
  University of California, San Diego\\
  \texttt{cow003@ucsd.edu} \\
}
\begin{document}
\maketitle


\begin{abstract}
The CSE 151B Spring 2026 competition tasks participants with maximising
accuracy on a mixed set of multiple-choice (MCQ) and free-form math problems
using a fixed open-weight reasoning model.
With no fine-tuning, all gains come from inference-time engineering.
We extend the starter notebook with five targeted changes:
(i)~BFloat16 full-precision serving instead of INT8 quantisation,
(ii)~scoped LaTeX repair inside \texttt{\$\ldots\$} delimiters only,
(iii)~structured six-step and five-step system prompts tuned for a thinking
model's \texttt{<think>/</think>} output format,
(iv)~a lazy-fallback MCQ extractor that invokes a cheap logprobs pass only
when the primary extractor abstains, and
(v)~brace-aware boxed-answer parsing.
These changes achieve a \textbf{final private leaderboard score of 0.615}
(\textbf{rank 37 out of 105 teams}, top 35\%).
Code: \url{https://github.com/Cjwu52104/151B_SP26_Competition}.
\end{abstract}


% % % === Introduction
\section{Introduction}
\label{sec:intro}

\paragraph{Problem Definition.}
The competition dataset contains $1{,}116$ problems partitioned into two
types: 375 multiple-choice questions (MCQ) each accompanied by up to ten
labelled options $A,B,\dots,J$, and 741 free-form problems whose gold answers
are symbolic or numeric expressions.
We are given a fixed model,
\texttt{Qwen3-4B-Thinking-2507}~\citep{qwen3}, and must produce an answer
for every question.
Predictions are submitted as a CSV; MCQ answers are scored by case-insensitive
letter equality and free-form answers by symbolic equivalence via a SymPy-based
judger.
No fine-tuning, retrieval, or external APIs are permitted at inference time.

\paragraph{Problem Significance.}
Automated mathematical reasoning is a core challenge for AI systems with
high-stakes applications in education, scientific computing, and software
verification.
Mid-size open-weight reasoning models running on a single consumer or
workstation GPU represent the most accessible deployment setting for students
and independent researchers.
Understanding which inference-time choices reliably improve accuracy —
particularly for ``thinking'' models that emit a long internal chain of thought
before their final answer — transfers directly to real low-compute
applications.

\paragraph{Technical Challenge.}
Three challenges shaped this project.
(1)~\textit{Corrupted input:} the dataset stores LaTeX with backslashes
stripped (\texttt{frac}, \texttt{infty}, \texttt{int}, etc.), and a naive
global repair corrupts common English words (``sum'', ``log'') that share
spellings with LaTeX commands.
(2)~\textit{Thinking-model output handling:} the model emits a long
\texttt{<think>}\ldots\texttt{</think>} internal monologue before its final
answer; the token budget and stop conditions must accommodate the full chain
without truncating it.
(3)~\textit{Robust answer extraction:} MCQ scoring requires a single letter
from a \texttt{\textbackslash boxed\{\}} expression, and free-form scoring
requires the response to contain exactly one well-formed
\texttt{\textbackslash boxed\{\}} whose content survives nested-brace parsing.

\paragraph{Contributions.}
The starter code is a working end-to-end pipeline (vLLM, 32k-token budget,
generic system prompts, regex letter extraction, symbolic judger).
Its limitations include:
no LaTeX repair,
INT8 quantisation that reduces effective precision,
MCQ extraction that never abstains (falls back to the last capital letter
anywhere in the response),
and truncation of nested \texttt{\textbackslash boxed\{\}} contents.
Our final solution addresses each limitation:

\begin{itemize}
  \item \textbf{BFloat16 full-precision serving.} Removing INT8 quantisation
    and serving in BF16 on a 24~GiB NVIDIA A30 eliminates weight rounding
    errors that can silently alter numerical answers.

  \item \textbf{Scoped LaTeX repair.} Backslashes are restored only inside
    \texttt{\$\ldots\$} math blocks using 35 pattern--replacement pairs,
    leaving prose untouched.

  \item \textbf{Structured thinking-tuned prompts.} The MCQ prompt directs the
    model through six explicit reasoning steps (Classify, Setup, Solve, Match,
    Eliminate, Confirm) and ends with a strict post-\texttt{</think>}
    instruction to emit exactly \texttt{\textbackslash boxed\{X\}}.
    The free-form prompt uses five steps (Count, Plan, Solve, Verify, Collect)
    and enforces a single comma-separated \texttt{\textbackslash boxed\{\}}.

  \item \textbf{Lazy-fallback MCQ extractor.} The thinking response is the
    primary source.  A cheap no-think logprobs pass over letter tokens $A$--$J$
    is invoked only when the primary extractor abstains, avoiding extra latency
    on the majority of questions.

  \item \textbf{Hardened answer parsing.} \texttt{extract\_letter} reads
    preferentially from the post-\texttt{</think>} region; letter-option
    markers (\texttt{A)}, \texttt{B.}) are used as a last resort.
\end{itemize}


% % % === Related Work
\section{Related Work}
\label{sec:relatedwork}

\paragraph{Reasoning LLMs.}
Recent open-weight ``thinking'' models such as the Qwen3
series~\citep{qwen3} interleave a hidden chain of thought delimited by
\texttt{<think>/</think>} with a final user-facing answer.
Their accuracy on math benchmarks improves with an adequate decoding budget
and sampling temperatures in the range recommended by the model authors
($T{=}0.6$, $p{=}0.95$, $k{=}20$).
Thinking models are sensitive to system-prompt wording: prompts that do
not explicitly instruct the model on post-think output format often elicit
verbose prose rather than a clean \texttt{\textbackslash boxed\{\}},
breaking downstream extraction.

\paragraph{Inference serving.}
vLLM~\citep{vllm} provides batched, paged-attention serving of HuggingFace
models with BitsAndBytes INT8 and BFloat16 precision modes.
Prefix caching (\texttt{enable\_prefix\_caching=True}) amortises the cost of
a shared system prompt across all questions in the batch, reducing
time-to-first-token for long prompts.

\paragraph{Symbolic answer judging.}
Math benchmarks such as MATH~\citep{hendrycks2021math} have driven the
development of symbolic equivalence checkers built on CAS backends.
The course-provided \texttt{Judger} class follows this lineage: it extracts
the content of \texttt{\textbackslash boxed\{\}} and tests mathematical
equivalence against each gold answer, tolerating equivalent forms such as
simplified fractions and expanded polynomials.


% % % === Methods
\section{Methods}
\label{sec:methods}

\paragraph{Problem Setting.}
Let $\mathcal{D}=\{(q_i,O_i,y_i)\}_{i=1}^{n}$ be the evaluation set, where
$q_i$ is the question text, $O_i$ is the (possibly empty) list of MCQ options,
and $y_i$ is the gold answer.
The model $f_\theta$ is held fixed.  We choose a prompt template $\pi$,
sampling parameters $\sigma$, and extraction function $g$ to maximise
\[
  \mathrm{Acc}(\pi,\sigma,g) \;=\; \tfrac{1}{n}\sum_{i=1}^{n}\,
    \mathbf{1}\!\left[\mathrm{Judge}\bigl(g(f_\theta(\pi(q_i,O_i);\sigma))\bigr) = y_i\right],
\]
where $\mathrm{Judge}$ is letter equality for MCQ and symbolic equivalence for
free-form.  Because $\theta$ is fixed, there is no trainable loss; all
optimisation is over the discrete choices $(\pi,\sigma,g)$.

\paragraph{Idea Summary.}
The core insight is that a thinking model already ``knows'' the answer —
its chain of thought is long and accurate — but failures arise at two
boundaries: (1)~the input boundary, where stripped LaTeX makes problem
statements ambiguous, and (2)~the output boundary, where imprecise prompts
and fragile extraction discard correct reasoning.
Targeting both boundaries with lightweight, zero-latency fixes is sufficient
to improve accuracy without any training or ensemble cost.
Alternative directions (self-consistency, SFT, GRPO) were queued but not
pursued, since the output-boundary fixes alone moved the leaderboard score
into the upper third.

\paragraph{Method Description.}
\Cref{tab:changes} summarises the differences between the starter code and
our final optimised pipeline.

\begin{table}[h]
  \caption{Component-level comparison of the starter and the final optimised
    pipeline. All changes are inference-time only; the model weights are
    unchanged.}
  \label{tab:changes}
  \centering
  \footnotesize
  \begin{tabular}{lp{4.3cm}p{4.3cm}}
    \toprule
    Component & Starter & Final (Optimised) \\
    \midrule
    Quantisation &
      INT8 via BitsAndBytes (\texttt{gpu\_memory\_utilization=0.50}). &
      BFloat16 full precision (\texttt{gpu\_memory\_utilization=0.90},
      NVIDIA A30 24~GiB). \\[3pt]
    Input preprocessing &
      None. &
      Scoped LaTeX repair: restores backslashes on 35 commands
      (\texttt{frac}, \texttt{sqrt}, \texttt{sin}, \texttt{pi}, \ldots)
      inside \texttt{\$\ldots\$} blocks only. \\[3pt]
    MCQ system prompt &
      Generic: ``output the letter in \texttt{\textbackslash boxed\{\}}.'' &
      Six-step structured: Classify $\to$ Setup $\to$ Solve $\to$ Match $\to$
      Eliminate $\to$ Confirm; strict post-\texttt{</think>} format rule. \\[3pt]
    Free-form system prompt &
      Generic step-by-step instruction. &
      Five-step structured: Count $\to$ Plan $\to$ Solve $\to$ Verify $\to$
      Collect; single comma-separated \texttt{\textbackslash boxed\{\}} rule. \\[3pt]
    MCQ extraction &
      Regex falls back to the last capital letter anywhere in the response
      (never abstains). &
      Multi-pattern extractor prefers text after \texttt{</think>};
      on abstention, a no-think logprobs pass reads $P(\text{letter})$
      over $A$--$J$. \\[3pt]
    Sampling / budget &
      $T{=}0.6$, $p{=}0.95$, $k{=}20$, \texttt{max\_tokens=32768},
      \texttt{enable\_prefix\_caching=False}. &
      Same temperature; \texttt{max\_tokens=16384};
      \texttt{enable\_prefix\_caching=True};
      \texttt{max\_num\_seqs=64}. \\
    \bottomrule
  \end{tabular}
\end{table}

\noindent
The \textbf{LaTeX repair} scans each question for \texttt{\$\ldots\$}
substrings and applies 35 pattern-replacement pairs inside those regions
only, preventing the repair from corrupting English words that share
spellings with LaTeX commands.

The \textbf{MCQ prompt} ends with: \textit{``After \texttt{</think>}, output
EXACTLY ONE LINE and nothing else: \texttt{\textbackslash boxed\{X\}} where
X is the single capital letter.''}
This eliminates the most common failure mode — prose explanation after the
think block that buries or omits the boxed answer.

The \textbf{free-form prompt} uses a COUNT step that explicitly labels each
\texttt{[ANS]} placeholder, ensuring multi-part answers are collected in
order into a single \texttt{\textbackslash boxed\{a\_1, a\_2, \ldots\}}.

The \textbf{lazy logprobs fallback} invokes a separate near-greedy pass
(\texttt{temperature=0.01}, \texttt{logprobs=20}, \texttt{max\_tokens=10})
with a minimal ``output only the answer letter'' prompt only for MCQs where
the primary extractor returns an empty string.
This keeps median latency flat while recovering answers on the small fraction
of questions where the main response is ambiguous.


% % % === Experiments
\section{Experiments}
\label{sec:experiments}

\subsection{Baselines}
\label{sec:baselines}

We compare against the following baselines:
\begin{itemize}
  \item \textbf{Starter (INT8)} -- the unmodified course-provided notebook,
    served with BitsAndBytes INT8 quantisation and generic prompts.
    Evaluated on the public development set as reported in milestone testing.

  \item \textbf{Optimised (no quant.)} -- our final submission, served in
    BF16 with structured prompts and robust extraction.
    This is the configuration submitted to the private leaderboard.
\end{itemize}

\subsection{Evaluation}
\label{sec:evaluation}

MCQ predictions are scored by case-insensitive letter equality with the gold
label.
Free-form predictions are scored by the course-provided \texttt{Judger},
which extracts the content of the first \texttt{\textbackslash boxed\{\}}
expression, normalises it, and tests mathematical equivalence against each
gold answer via SymPy-based symbolic evaluation.
Overall accuracy is reported as total correct / total questions, with MCQ
and free-form items weighted equally.
The final private score comes from the Kaggle leaderboard.

\subsection{Implementation Details}
\label{sec:implementation_details}

\textbf{Hardware.}
All inference runs on a single NVIDIA A30 GPU (24~GiB VRAM).
A full 1,116-question run takes approximately 4--6 hours, which makes
brute-force hyperparameter sweeps impractical and motivates our focus on
design choices with high expected return per unit of reasoning.

\textbf{Model and serving.}
The model is \texttt{Qwen/Qwen3-4B-Thinking-2507} served by vLLM~\citep{vllm}
with \texttt{dtype="bfloat16"}, \texttt{max\_model\_len=16384},
\texttt{gpu\_memory\_utilization=0.90}, \texttt{enable\_prefix\_caching=True},
and \texttt{max\_num\_seqs=64}.
We removed BitsAndBytes quantisation relative to the starter because the A30
has sufficient VRAM to hold the 4B-parameter model in BF16
($\approx$8~GiB weights), leaving headroom for a large KV cache.

\textbf{Sampling.}
We follow the Qwen3-Thinking recommended defaults:
$T{=}0.6$, $p{=}0.95$, $\mathrm{top\_k}{=}20$,
$\mathrm{repetition\_penalty}{=}1.0$, \texttt{max\_tokens=16384}.
For the logprobs calibration fallback, we use $T{=}0.01$, $p{=}1.0$,
\texttt{logprobs=20}, \texttt{max\_tokens=10}.

\textbf{Prompts.}
The system prompts are applied via the tokenizer's
\texttt{apply\_chat\_template} with \texttt{add\_generation\_prompt=True};
no few-shot examples are included.
The MCQ prompt is 109 tokens and the free-form prompt is 127 tokens; both
benefit from prefix caching across the batch.

\textbf{Post-processing.}
For MCQs, \texttt{extract\_letter} searches the post-\texttt{</think>}
region first using five regex patterns (boxed letter, ``answer is X'',
``answer: X'', ``therefore the answer is X'', ``option X is correct'')
before falling back to option-marker heuristics.
Questions that return an empty string trigger the logprobs calibration pass.
For free-form, the raw response is passed directly to the \texttt{Judger}.

\subsection{Results}
\label{sec:results}

\paragraph{Final leaderboard.}
\Cref{tab:final-results} shows the key result: our optimised pipeline
achieves \textbf{0.615} overall accuracy on the private test set, placing
\textbf{37th out of 105 teams} (top 35\%).

\begin{table}[h]
  \caption{Final private leaderboard result.  The starter baseline is
    approximated from small-$N$ public-set testing; only the optimised score
    is from the official leaderboard.}
  \label{tab:final-results}
  \centering
  \begin{tabular}{lcc}
    \toprule
    Configuration & Overall Accuracy & Leaderboard Rank \\
    \midrule
    Starter (INT8, generic prompts, $N{\approx}5$) & $\approx$0.60 & -- \\
    \textbf{Optimised (BF16, structured prompts)} & \textbf{0.615} & \textbf{37 / 105} \\
    \bottomrule
  \end{tabular}
\end{table}

\paragraph{Qualitative findings.}
Manual inspection of model responses reveals two recurring failure patterns.
First, about 5--10\% of free-form responses emit two or more
\texttt{\textbackslash boxed\{\}} expressions — one inside the think block
and one after it — causing the judger to read the inner (partial) boxed
expression.
Second, for multi-part questions with three or more sub-answers, the model
occasionally switches ordering mid-response, producing a boxed tuple whose
entries are correct but in the wrong sequence.
Both patterns suggest that further prompt engineering — explicitly warning
against boxed content inside \texttt{<think>} and repeating the enumeration
order — could recover additional accuracy.


\subsection{Ablations}
\label{sec:ablations}

Due to the 4--6 hour per-run cost on our single GPU, we did not run a full
ablation sweep over every component in isolation.
Instead, we developed on a frozen 100-item stratified subsample
(75 free-form, 25 MCQ) from the public development set to test each change
before committing to a full run.
\Cref{tab:ablations} summarises the expected and observed contributions of
each component.

\begin{table}[h]
  \caption{Ablation summary. ``Expected direction'' is based on failure-mode
    analysis on the 100-item subsample.  ``Latency impact'' indicates
    whether the change adds inference cost.}
  \label{tab:ablations}
  \centering
  \footnotesize
  \begin{tabular}{llll}
    \toprule
    Component & Expected direction & Latency impact & Rationale \\
    \midrule
    BF16 vs.\ INT8 &
      $+$ accuracy &
      Neutral (faster per token) &
      Eliminates weight-rounding on precise numeric answers \\
    Scoped LaTeX repair &
      $+$ MCQ and free-form &
      None &
      Prevents garbled math notation from misleading the model \\
    Structured MCQ prompt &
      $+$ MCQ &
      None (same token count) &
      Forces single clean \texttt{\textbackslash boxed\{X\}} after \texttt{</think>} \\
    Structured free-form prompt &
      $+$ free-form &
      None &
      Forces ordered comma-separated sub-answers into one \texttt{\textbackslash boxed\{\}} \\
    Lazy logprobs fallback &
      $+$ MCQ (abstained) &
      Marginal ($<$5\% of MCQs) &
      Recovers letter when main extractor returns empty string \\
    \texttt{enable\_prefix\_caching} &
      $\pm$ (throughput) &
      Reduces TTFT &
      Amortises 100+ token shared system prompt across batch \\
    \bottomrule
  \end{tabular}
\end{table}

Of these, the switch from INT8 to BF16 and the structured MCQ prompt are the
two changes with the highest expected return: INT8 quantisation introduces
rounding errors that compound in multi-step arithmetic, and a poorly-formatted
MCQ response fails extraction entirely rather than partially.
Both changes were validated on the 100-item subsample before the final full run.


% % % === Discussion
\section{Discussion}
\label{sec:discussion}

Thinking models present a subtle deployment challenge: they are unusually
sensitive to anything that prevents the final answer from surfacing.
Token budgets that truncate the chain, extraction paths that prefer a
no-think pass, and \texttt{stop=[\ldots]} lists that fire during reasoning all
act as silent answer-suppressors.
Our work addresses these failure modes at both the input and output boundaries
without requiring any training or ensemble cost.

\paragraph{Achievements.}
Our inference-only pipeline achieves 0.615 overall accuracy on the private
test set, ranking 37th out of 105 teams.
All gains come from five targeted engineering decisions — BF16 serving,
scoped LaTeX repair, structured thinking-tuned prompts, lazy logprobs
fallback, and prefix caching — each individually motivated by a specific
observed failure mode.
The final pipeline is a single Python script (\texttt{run\_inference.py})
that reproduces the submitted result in one command with no external
dependencies beyond \texttt{vllm} and \texttt{transformers}.

\paragraph{Bottlenecks and Mitigation.}
The dominant bottleneck was decode latency: the 4--6 hour full-dataset
runtime made iterative hyperparameter sweeps infeasible.
We mitigated this by maintaining a 100-item stratified subsample as a fast
evaluation proxy.
A secondary bottleneck was the multi-boxed failure mode noted above — the
judger reads the \emph{first} \texttt{\textbackslash boxed\{\}} expression,
which on roughly 5--10\% of free-form responses is an intermediate result
inside the think block rather than the final answer.
A post-processing step that skips boxed content inside
\texttt{<think>}\ldots\texttt{</think>} would fix this without any model
change, but was not implemented before the submission deadline.

\paragraph{Lessons Learned.}
The most important lesson is that output-boundary failures are both common
and easily fixable for thinking models.
Generic prompts that do not anticipate the \texttt{<think>/</think>} format
produce responses where the answer is technically present but
extraction fails.
Structured prompts with explicit post-think instructions require no extra
computation and eliminate this failure class almost entirely.

A second lesson is that precision matters for math: INT8 quantisation
introduces compounding rounding errors on multi-step arithmetic that a
higher-precision format avoids.
Given sufficient VRAM, BF16 serving should be preferred for math-reasoning
tasks even at the cost of slightly higher memory usage.

Finally, a frozen evaluation subsample is essential for fast iteration on a
budget of single-GPU hours.
Without it, each change requires a full multi-hour run, making engineering
feedback loops far too slow.

\paragraph{Team Responsibilities.}
Solo submission.
Connor Wu designed and implemented the full pipeline, including LaTeX repair,
system prompts, the logprobs fallback, the inference harness, and the
evaluation setup.
All experiments were run and all results were verified by the author.


\section*{LLM Usage}

Anthropic's Claude (model \texttt{claude-sonnet-4-6}, via the Claude Code
desktop interface) was used as an interactive coding assistant with
read/write access to this repository.
Claude proposed and applied code changes, and drafted portions of this report
from briefings on scope, authorship, and results.
The author designed the overall approach, ran all experiments, verified all
code changes, and reviewed and revised every factual claim in the report.
The author takes full responsibility for the contents.


\bibliography{references.bib}
\bibliographystyle{plainnat}

\end{document}
